% Diapositivas de avance (datos y procesamiento). Compilar desde esta carpeta:  latexmk -pdf diapositivas.tex
% Las figuras se leen de ../resultados/paso2/ (generadas por paso2_features.py y material_reporte.py).
\documentclass[aspectratio=169,11pt]{beamer}
\usetheme[progressbar=frametitle, numbering=fraction, sectionpage=none, block=fill]{metropolis}
\usepackage[T1]{fontenc}
\usepackage[spanish,es-noshorthands,es-tabla]{babel}
\usepackage{booktabs}
\usepackage{array}
\newcolumntype{P}[1]{>{\raggedright\arraybackslash}p{#1}}   % columna de texto alineada a la izquierda
\usepackage{tikz}
\usetikzlibrary{positioning,arrows.meta}
\usepackage{siunitx}
\sisetup{output-decimal-marker={,}}
\graphicspath{{../resultados/paso2/}}

\definecolor{azul}{HTML}{2A78D6}
\definecolor{gris}{HTML}{52514E}
\setbeamercolor{progress bar}{fg=azul, bg=azul!15}
\setbeamercolor{alerted text}{fg=azul}
\setbeamercolor{block title}{bg=azul!18, fg=black}
\setbeamercolor{block body}{bg=azul!6, fg=black}
\setbeamercolor{block title alerted}{bg=azul, fg=white}
\setbeamercolor{block body alerted}{bg=azul!6, fg=black}
\newcommand{\uV}{\si{\micro\volt}}
\makeatletter  % \input primitivo: filas de tablas generadas por material_reporte.py desde los CSV
\newcommand{\tabla}[1]{\@@input ../resultados/paso2/tablas/#1 }
\makeatother

\title{Características EEG: ventana y solapamiento}
\subtitle{TEPT, Alzheimer y depresión mayor. Avance: datos y procesamiento}
\author{Aldrin Chavez, Jhony Peñaherrera, Marck Chiza y Fernando Chantera}
\date{24 de noviembre de 2026}

\begin{document}

\maketitle

% ---------------------------------------------------------------------------------------------
\begin{frame}{Punto de partida: OnTEPT}
  \begin{itemize}
    \item OnTEPT (Salazar-Castro et al., 2026) es una ontología de TEPT que integra \alert{36 características EEG} de 41 participantes.
    \item No documenta \alert{con qué parámetros} se calcularon (ventana, solapamiento) ni \alert{qué tan fiables} son.
  \end{itemize}
  \vspace{4mm}
  \begin{block}{Pregunta}
    ¿Cómo cambian la \textbf{estabilidad}, el \textbf{costo} y la \textbf{relevancia} (SHAP) de las características EEG según la longitud de ventana (1, 2 y 4 s) y el solapamiento (0, 50 y 75 \%)?
  \end{block}
  \vspace{3mm}
  {\small Análisis exploratorio de investigación, no clínico: TEPT y otras dos enfermedades.}
\end{frame}

% ---------------------------------------------------------------------------------------------
\begin{frame}{Plan en tres pasos}
  \begin{enumerate}
    \item \textbf{Datos}: tres conjuntos públicos, uno de TEPT y dos de otras enfermedades. (\alert{hecho})
    \item \textbf{Nueve configuraciones}: ventanas de 1, 2 y 4 s por solapamientos de 0, 50 y 75 \%; se miden la estabilidad y el tiempo. (\alert{hecho})
    \item \textbf{Clasificador y SHAP} en cada conjunto: qué características, bandas y canales pesan más. (pendiente)
  \end{enumerate}
  \vspace{5mm}
  \begin{alertblock}{Regla}
    Nunca mezclar conjuntos en un mismo modelo: cada uno usa un equipo distinto, y el modelo aprendería a reconocer el equipo y no la enfermedad.
  \end{alertblock}
\end{frame}

% ---------------------------------------------------------------------------------------------
\begin{frame}{Búsqueda de datos}
  \begin{columns}[T]
    \begin{column}{0.36\textwidth}
      \textbf{Criterios}
      \begin{itemize}\small
        \item Acceso abierto
        \item Diagnóstico etiquetado
        \item Controles sanos
        \item Registro en reposo
        \item Canales y frecuencia documentados
      \end{itemize}
    \end{column}
    \begin{column}{0.6\textwidth}
      \textbf{Fuentes de EEG citadas por OnTEPT}
      \begin{itemize}\small
        \item du Bois 2022: 29 pacientes con TEPT, \alert{sin controles sanos}; EEG solo en 20.
        \item Rahmani 2018: 6 con TEPT y 6 controles. \alert{El DOI de datos que da OnTEPT es erróneo.}
        \item DAIC-WOZ: solo cuestionarios, sin EEG.
      \end{itemize}
      \textbf{Otras fuentes}
      \begin{itemize}\small
        \item OpenNeuro y TDBRAIN: ningún EEG de TEPT.
        \item Park 2021: solo características ya calculadas, sin señal.
      \end{itemize}
    \end{column}
  \end{columns}
\end{frame}

% ---------------------------------------------------------------------------------------------
\begin{frame}{Los tres conjuntos de datos}
  \footnotesize
  \centering
  \begin{tabular}{@{}P{2.6cm} P{3.3cm} P{3.3cm} P{3.3cm}@{}}
    \toprule
     & \textbf{TEPT} & \textbf{Alzheimer} & \textbf{Depresión mayor} \\
    \midrule
    Fuente & Rahmani 2018 (figshare) & ds004504 (OpenNeuro) & Mumtaz 2017 (figshare) \\
    Pacientes / controles & 6 / 6 & 36 / 29 & 28 / 27 \\
    Registro & Dentro de un resonador (EEG con fMRI) & Clínico, reposo, ojos cerrados & Reposo, ojos cerrados \\
    Canales & 31 + ECG & 19 & 19 + auxiliares \\
    Muestreo & 250 Hz & 500 Hz & 256 Hz \\
    Referencia & FCz & Cz & Orejas unidas \\
    Duración & 520 s & 5--21 min & $\approx$ 5 min \\
    Procesado previo & Sí (ICA, filtro a 40 Hz) & No & No \\
    \bottomrule
  \end{tabular}

  \vspace{3mm}
  {\small Todo difiere entre conjuntos: hay que \alert{armonizar} la señal y entrenar \alert{un modelo por conjunto}.}
\end{frame}

% ---------------------------------------------------------------------------------------------
\begin{frame}{Control de calidad}
  \begin{columns}[T]
    \begin{column}{0.6\textwidth}
      \small
      \begin{tabular}{@{}P{4.1cm} P{3.7cm}@{}}
        \toprule
        Problema en depresión & Qué se hizo \\
        \midrule
        6 archivos rotos en figshare & No se pueden descargar \\
        2 pares de archivos idénticos & Se elimina la copia \\
        1 registro de 180 s & Se excluye ($<$ 240 s) \\
        Canales extra casi solo en controles & Se usan los 19 comunes \\
        \bottomrule
      \end{tabular}

      \vspace{3mm}
      {\footnotesize Con un duplicado en entrenamiento y otro en prueba, el modelo reconocería la señal y no la enfermedad: la exactitud saldría inflada.}
    \end{column}
    \begin{column}{0.36\textwidth}
      \begin{block}{Muestra final: 132 sujetos}
        TEPT: 6 + 6\\
        Alzheimer: 36 + 29\\
        Depresión: 28 + 27\\[2pt]
        {\footnotesize (pacientes + controles)}
      \end{block}
    \end{column}
  \end{columns}
\end{frame}

% ---------------------------------------------------------------------------------------------
\begin{frame}{El mismo procesamiento para los tres conjuntos}
  \centering
  \begin{tikzpicture}[
      caja/.style={draw=azul, fill=azul!8, rounded corners=2pt, align=center, font=\scriptsize,
                   minimum height=12mm, text width=19mm, inner sep=2pt},
      analisis/.style={caja, draw=gris, fill=gris!8},
      flecha/.style={-{Stealth[length=2mm]}, semithick, draw=gris},
      node distance=10mm and 3mm]
    \node[caja] (a) {Lectura:\\19 canales en \uV};
    \node[caja, right=of a] (b) {Remuestreo\\a 256 Hz};
    \node[caja, right=of b] (c) {Filtro\\1--40 Hz};
    \node[caja, right=of c] (d) {Referencia\\promedio};
    \node[caja, right=of d] (e) {Recorte\\a 240 s};
    \node[caja, right=of e] (f) {Segundos con\\artefactos};
    \node[analisis, below=of f] (g) {Ventanas:\\9 configuraciones};
    \node[analisis, left=of g] (h) {228 caracte-\\rísticas};
    \node[analisis, left=of h] (i) {Estabilidad\\y tiempo};
    \node[analisis, left=of i, dashed] (j) {Clasificador\\y SHAP\\(pendiente)};
    \foreach \x/\y in {a/b, b/c, c/d, d/e, e/f, f/g, g/h, h/i, i/j}
      \draw[flecha] (\x) -- (\y);
  \end{tikzpicture}

  \vspace{5mm}
  {\small \textcolor{azul}{Arriba}: armonizar la señal. \textcolor{gris}{Abajo}: cortar en ventanas, calcular características y evaluar.}
\end{frame}

% ---------------------------------------------------------------------------------------------
\begin{frame}{Armonización: qué se hace y por qué}
  \footnotesize
  \centering
  \begin{tabular}{@{}l P{4.2cm} P{7.4cm}@{}}
    \toprule
    Paso & Qué & Por qué \\
    \midrule
    Canales & 19 del sistema 10-20 & Son los únicos comunes a los tres conjuntos \\
    Muestreo & 256 Hz & La entropía depende de la frecuencia; con 256 Hz los pasos de ventana son enteros \\
    Filtro & 1--40 Hz, fase cero & Quita derivas lentas; 40 Hz es el techo de Rahmani; elimina la red eléctrica \\
    Referencia & Promedio de los 19 canales & Cada conjunto usaba otra; la original se cancela \\
    Duración & 240 s por sujeto & Misma cantidad de datos y mismo peso por sujeto \\
    Artefactos & Segundo malo si supera \SI{150}{\micro\volt} pico a pico & Quita parpadeos y movimientos sin perder el ritmo alfa amplio \\
    \bottomrule
  \end{tabular}

  \vspace{2mm}
  {\small La \alert{estandarización de las características} (z-score) irá en el paso 3, dentro de la validación cruzada.}
\end{frame}

% ---------------------------------------------------------------------------------------------
\begin{frame}{Qué hace el procesamiento: un ejemplo real}
  \centering
  \includegraphics[width=0.96\textwidth]{fig_ejemplo_senal.png}

  {\footnotesize Control sub-050 (ds004504), canal O1. A: se van la deriva y el ruido. B: una ventana sola es ruidosa; el promedio, estable.}
\end{frame}

% ---------------------------------------------------------------------------------------------
\begin{frame}{Ventanas y solapamiento}
  \begin{columns}[c]
    \begin{column}{0.5\textwidth}
      \centering
      \begin{tikzpicture}[x=6mm, y=1mm]
        \foreach \base/\paso/\etq in {0/4/{0 \%}, -12/2/{50 \%}, -28/1/{75 \%}} {
          \node[anchor=east, font=\footnotesize] at (-0.2, \base-1) {\etq};
          \pgfmathtruncatemacro{\n}{4/\paso}
          \foreach \k in {0,...,\n} {
            \pgfmathsetmacro{\dy}{(\paso<4)*\k*2.6}
            \fill[azul!20, draw=azul] (\k*\paso, \base-\dy) rectangle ++(4, -2);
          }
        }
        \draw[-{Stealth}, gris] (0, -44) -- (8.6, -44) node[right, font=\scriptsize] {s};
        \foreach \t in {0, 2, 4, 6, 8}
          \draw[gris] (\t, -43.3) -- (\t, -44.7) node[below, font=\scriptsize] {\t};
      \end{tikzpicture}

      {\scriptsize Ventanas de 4 s durante los primeros 8 s}
    \end{column}
    \begin{column}{0.48\textwidth}
      \footnotesize
      \centering
      \begin{tabular}{@{}lrrr@{}}
        \toprule
         & \multicolumn{3}{c}{Ventanas en 240 s} \\
        \cmidrule(l){2-4}
        Ventana & 0 \% & 50 \% & 75 \% \\
        \midrule
        1 s & 240 & 479 & 957 \\
        2 s & 120 & 239 & 477 \\
        4 s & 60 & 119 & 237 \\
        \bottomrule
      \end{tabular}

      \vspace{3mm}
      \begin{itemize}
        \item Más solapamiento: más ventanas, pero \alert{no más información}.
        \item Ventana más larga: mejor resolución en frecuencia (de 1 a 0,25 Hz), pero más pérdida por artefactos.
      \end{itemize}
    \end{column}
  \end{columns}
\end{frame}

% ---------------------------------------------------------------------------------------------
\begin{frame}{228 características por ventana}
  \footnotesize
  \centering
  \begin{tabular}{@{}P{3.9cm} c P{5.4cm} >{\centering\arraybackslash}p{2.3cm}@{}}
    \toprule
    Familia & Tipos & Qué mide & ¿Depende de la amplitud? \\
    \midrule
    Potencia absoluta ($\log_{10}$) & 5 & Cuánta actividad hay en la banda & Sí \\
    Potencia relativa & 5 & Proporción del total de 1--40 Hz & No \\
    Entropía espectral & 1 & Qué tan repartido está el espectro (0 a 1) & No \\
    Entropía de permutación & 1 & Qué tan impredecible es la forma de onda (0 a 1) & No \\
    \bottomrule
  \end{tabular}

  \vspace{4mm}
  \raggedright
  \begin{itemize}\small
    \item Bandas: $\delta$ 1--4, $\theta$ 4--8, $\alpha$ 8--13, $\beta$ 13--30 y $\gamma$ 30--40 Hz.
    \item 12 tipos × 19 canales = \alert{228 características}; por ejemplo, \texttt{rel\_alpha\_O1}.
    \item Espectro: periodograma con ventana de Hann; su resolución es 1 / longitud de la ventana.
  \end{itemize}
\end{frame}

% ---------------------------------------------------------------------------------------------
\begin{frame}{¿Cómo se mide la estabilidad?}
  \centering
  \begin{tikzpicture}[font=\footnotesize]
    \draw[fill=azul!15, draw=azul] (0, 0) rectangle (5, 0.7) node[midway] {1.ª mitad: 0--120 s};
    \draw[fill=azul!35, draw=azul] (5, 0) rectangle (10, 0.7) node[midway] {2.ª mitad: 120--240 s};
    \node[anchor=east] at (-0.1, 0.35) {Registro};
    \draw[-{Stealth}, gris] (2.5, -0.1) -- (2.5, -0.7) node[below] {promedio A};
    \draw[-{Stealth}, gris] (7.5, -0.1) -- (7.5, -0.7) node[below] {promedio B};
    \node[draw=gris, rounded corners=2pt, fill=gris!8, align=center, inner sep=4pt] at (5, -2.3)
      {ICC(A,1) entre sujetos: ¿cada sujeto conserva su valor de A a B?};
  \end{tikzpicture}

  \vspace{5mm}
  \begin{columns}[T]
    \begin{column}{0.5\textwidth}
      \footnotesize
      \begin{itemize}
        \item ICC $\approx$ 1: la característica «reconoce» al sujeto.
        \item ICC $\approx$ 0: cambia entre mitades tanto como entre personas.
      \end{itemize}
    \end{column}
    \begin{column}{0.44\textwidth}
      \footnotesize
      Escala de Koo y Li (2016): menos de 0,5, pobre; 0,5--0,75, moderada; 0,75--0,9, buena; más de 0,9, excelente.

      \vspace{2mm}
      Mismos sujetos en las 9 configuraciones.
    \end{column}
  \end{columns}
\end{frame}

% ---------------------------------------------------------------------------------------------
\begin{frame}{Resultado 1: la estabilidad es alta en todas las configuraciones}
  \centering\small
  ICC mediano entre mitades: mínimo y máximo en las 9 configuraciones\\[2mm]
  \begin{tabular}{@{}lccc@{}}
    \toprule
    Familia & TEPT ($n = 12$) & Alzheimer ($n = 65$) & Depresión ($n = 55$) \\
    \midrule
    \tabla{estabilidad_familias.tex}
    \bottomrule
  \end{tabular}

  \vspace{4mm}
  \raggedright
  \begin{itemize}\small
    \item El solapamiento no cambia la estabilidad y la longitud de ventana, casi tampoco.
    \item Menos estables: $\delta$ y $\alpha$ relativas en TEPT (ICC de 0,6 a 0,8) y $\gamma$ relativa en Alzheimer y depresión.
  \end{itemize}
\end{frame}

% ---------------------------------------------------------------------------------------------
\begin{frame}{Resultado 2: el solapamiento solo añade costo}
  \centering\small
  \begin{tabular}{@{}lcccc@{}}
    \toprule
     & \multicolumn{3}{c}{Tiempo por sujeto (ms)} & Ventanas descartadas \\
    \cmidrule(lr){2-4}
    Ventana & 0 \% & 50 \% & 75 \% & por artefactos (\%) \\
    \midrule
    \tabla{costo_diapositivas.tex}
    \bottomrule
  \end{tabular}

  {\footnotesize Medianas de los 132 sujetos (240 s y 19 canales cada uno). Descartadas: rango entre los tres conjuntos, sin solapamiento.}

  \vspace{1mm}
  \raggedright
  \begin{itemize}\small
    \item Con 50 \% el tiempo se duplica y con 75 \% se multiplica por unas 7, con la misma estabilidad.
    \item Las ventanas largas pierden más datos: basta un segundo malo para descartar la ventana entera.
  \end{itemize}
\end{frame}

% ---------------------------------------------------------------------------------------------
\begin{frame}{Para tener en cuenta}
  \begin{itemize}
    \item \textbf{Amplitud}: en TEPT y depresión los pacientes tienen cerca de la mitad de amplitud que los controles (RMS de 4,0 frente a 7,8 \uV\ y de 5,3 frente a 8,0 \uV). La potencia absoluta podría separar los grupos por el registro y no por la enfermedad.
    \item \textbf{TEPT}: solo 12 sujetos, registrados dentro de un resonador y ya preprocesados.
    \item \textbf{Techo del ICC}: con 2 minutos por mitad, los promedios son muy estables; las diferencias entre configuraciones deberían verse ventana a ventana.
    \item \textbf{Calidad de datos}: el conjunto de depresión perdió 9 de sus 64 registros.
  \end{itemize}
\end{frame}

% ---------------------------------------------------------------------------------------------
\begin{frame}{Próximo paso: clasificador y SHAP}
  \begin{itemize}
    \item Un clasificador simple por conjunto (pacientes frente a controles) para cada una de las 9 configuraciones.
    \item Validación cruzada \alert{por sujeto}: dejando uno fuera en TEPT y con 5 particiones en Alzheimer y depresión.
    \item Estandarización dentro de cada partición de entrenamiento.
    \item SHAP: qué características, bandas y canales pesan más, y si cambian con la configuración.
  \end{itemize}
  \vspace{2mm}
  \begin{block}{Por decidir}
    El tipo de clasificador (propuesta: regresión logística con L2) y si se compara un modelo con todas las características frente a otro sin potencia absoluta.
  \end{block}
\end{frame}

% ---------------------------------------------------------------------------------------------
\begin{frame}{Referencias principales}
  \scriptsize
  \begin{itemize}
    \item Salazar-Castro, J. A. et al. (2026). OnTEPT: a multimodal OWL ontology for post-traumatic stress disorder. \emph{Journal of Biomedical Semantics}, 17:21.
    \item Rahmani, B. et al. (2018). Dynamical Hurst analysis identifies EEG channel differences between PTSD and healthy controls. \emph{PLOS ONE}, 13(7):e0199144. Datos: doi 10.6084/m9.figshare.5393338.
    \item Miltiadous, A. et al. (2023). A dataset of scalp EEG recordings of Alzheimer's disease, frontotemporal dementia and healthy subjects from routine EEG. \emph{Data}, 8(6):95. Datos: OpenNeuro ds004504.
    \item Mumtaz, W. et al. (2017). A wavelet-based technique to predict treatment outcome for major depressive disorder. \emph{PLOS ONE}, 12(2):e0171409. Datos: doi 10.6084/m9.figshare.4244171.
    \item du Bois, N. et al. (2022). \emph{Data in Brief}, 42:108066.
    \item McGraw, K. O. y Wong, S. P. (1996). \emph{Psychological Methods}, 1(1):30--46.
    \item Koo, T. K. y Li, M. Y. (2016). \emph{Journal of Chiropractic Medicine}, 15(2):155--163.
    \item Bandt, C. y Pompe, B. (2002). \emph{Physical Review Letters}, 88:174102.
    \item Inouye, T. et al. (1991). \emph{Electroencephalography and Clinical Neurophysiology}, 79(3):204--210.
  \end{itemize}
  \vspace{1mm}
  {\footnotesize Lista completa con DOI en el informe (carpeta \texttt{Report}).}
\end{frame}

\end{document}
